\section{Experiments}
\label{sec:experiments}

\begin{table}[t]
\centering\footnotesize
\setlength{\tabcolsep}{1.7pt}
\begin{tabular}{lcc|ccccc|cc}
\toprule
& \multicolumn{2}{c|}{Training set} & \multicolumn{5}{c|}{Per-class accuracy} & & \\
Model & \#Ros. & Noise & Dai. & Dan. & Ros. & Sun. & Tul. & Acc. & F1 \\
\midrule
base & 200 & 0 & 92 & 84 & 84 & 91 & 82 & 86.6 & 86.6 \\
\midrule
imb\_roses50 & 50 & 0 & 90 & 86 & 59 & 91 & 89 & 83.0 & 82.7 \\
imb\_roses20 & 20 & 0 & 92 & 87 & 36 & 91 & 91 & 79.4 & 78.3 \\
imb\_roses5 & 5 & 0 & 93 & 86 & \textbf{2} & 92 & 90 & 72.6 & 66.5 \\
\midrule
noise10 & 200 & 10 & 83 & 79 & 80 & 87 & 79 & 81.6 & 81.6 \\
noise30 & 200 & 30 & 75 & 72 & 70 & 76 & 67 & 72.0 & 72.0 \\
noise50 & 200 & 50 & 54 & 53 & 56 & 60 & 48 & 54.2 & 54.2 \\
\bottomrule
\end{tabular}
\caption{Training sets and test results (\%). \#Ros.: roses training images (others: 200); Noise: \% of training labels flipped. Per-class accuracy is recall; F1 is macro-F1.}
\label{tab:results}
\end{table}


The balanced \texttt{base} model reaches 86.6\%; its most common mistake is confusing roses and tulips
(9 images each way), which share red and pink petals.

\paragraph{RQ1: class imbalance.}
As roses training images drop from 200 to 50, 20, and 5, roses accuracy falls from 84\% to 59, 36, and 2\%.
Even a 4:1 imbalance costs 25 points.
The other four classes, in contrast, hardly change (82--93\%), and tulips actually improves from 82 to 90\%.
\Cref{fig:main}c explains why: the roses the model no longer recognizes are mostly labeled as tulips
(9, 29, 56, and 80 of 100), the class they already resembled.
With only 5 roses, the model predicts tulips 184 times for 100 real tulips, so tulip precision drops to 49\%.
Overall accuracy (72.6\%) understates this failure, since roses is only 20\% of the test set;
macro-F1 (66.5\%) reveals it more clearly.

\paragraph{RQ2: label noise.}
With 10, 30, and 50\% of the labels flipped, accuracy falls from 86.6\% to 81.6, 72.0, and 54.2\%:
about 0.5 points per 1\% of noise up to 30\%, then 0.9 points per 1\% beyond it.
Unlike imbalance, noise hurts all classes about equally: the per-class band in \cref{fig:main}b stays narrow,
and macro-F1 stays within 0.1 points of accuracy.
The extra errors still fall mostly between similar classes; tulips predicted as roses rise from 9 to 18 and 23
at 0, 30, and 50\% noise.
Even at 50\% noise, the correct label remains the most common one in each class (50\% versus 12.5\% for each
wrong class), which keeps the model well above chance (20\%).

\paragraph{Comparison and limitations.}
Flipping 100 labels (\texttt{noise10}) costs 5.0 points across all classes, whereas removing 150 roses
(\texttt{imb\_roses50}) costs 3.6 points overall but 25 on roses.
Each model was trained once, and with 100 test images per class the 95\% interval is about $\pm$7 points per class
($\pm$3 overall), so small gaps such as tulips rising from 82 to 90\% are only suggestive.
We also tested a single minority class and symmetric synthetic noise, whereas real label errors tend to cluster
between similar classes.
